\documentclass[10pt]{article}
\usepackage[margin=1in]{geometry}
\usepackage[T1]{fontenc}
\usepackage{lmodern,microtype,amsmath,amssymb,amsthm}
\theoremstyle{definition}
\newtheorem{assumption}{Assumption}
\begin{document}
\section*{Generalized Lipschitz Minty Root Finding}
We seek a zero of a continuous operator $F:\mathbb{R}^d\to\mathbb{R}^d$ in Euclidean dimension $d\geq 0$. On each segment, $F$ is Lipschitz with constant $L_0+L_1 M$, where $L_0>0$, $L_1\geq 0$, and $M$ is the maximum of $\|F\|$ on that segment. A given starting point $x^0$ lies within distance $R\geq 0$ of a zero $x^\star$ satisfying $\langle F(x),x-x^\star\rangle\geq 0$ for every $x\in\mathbb{R}^d$. The oracle returns $F(x)$ at $x$. For $\varepsilon>0$, the goal is an output $\bar{x}\in\mathbb{R}^d$ with $\mathbb{E}\|F(\bar{x})\|\leq\varepsilon$, where the expectation is over the algorithm's randomness.

In terms of the Euclidean residual norm, the problem is
\[
\min_{x\in\mathbb{R}^d}\|F(x)\|.
\]
The algorithm receives $d\in\{0,1,2,\ldots\}$, $L_0$, $L_1$, $R$, $\varepsilon$, and $x^0\in\mathbb{R}^d$, and may query any point in $\mathbb{R}^d$.

Our assumptions:
\begin{assumption}[Generalized Lipschitz continuity]
The operator $F$ is continuous. For every $x,y\in\mathbb{R}^d$, let $M_{x,y}$ denote the maximum operator value in norm along the segment from $x$ to $y$. Then
\[
\begin{aligned}
M_{x,y} &= \max_{t\in[0,1]}\|F(x+t(y-x))\|,\\
\|F(x)-F(y)\| &\leq (L_0+L_1 M_{x,y})\|x-y\|.
\end{aligned}
\]
\end{assumption}

\begin{assumption}[Minty zero and initial distance]
There exists $x^\star\in\mathbb{R}^d$ such that $F(x^\star)=0$, $\|x^0-x^\star\|\leq R$, and
\[
\langle F(x),x-x^\star\rangle\geq 0
\qquad\text{for every }x\in\mathbb{R}^d.
\]
\end{assumption}

After the allotted oracle calls, the residual norm $\|F(\bar{x})\|$ must be integrable under the law $P$ of the algorithm's randomness and satisfy
\[
\mathbb{E}_{P}\bigl[\|F(\bar{x})\|\bigr]\leq\varepsilon.
\]
\end{document}
